\section{Giới thiệu}
\label{sec:introduction}

Khai thác tập phổ biến (Frequent Itemset Mining, FIM) là một bài toán nền tảng trong khai thác dữ liệu giao dịch. Với một cơ sở dữ liệu gồm nhiều giao dịch, mục tiêu của FIM là tìm các nhóm item xuất hiện cùng nhau đủ thường xuyên theo một ngưỡng hỗ trợ tối thiểu. Kết quả của bài toán này thường được dùng làm đầu vào cho luật kết hợp, phân tích giỏ hàng, phát hiện mẫu trong log hệ thống, hoặc phân tích các tổ hợp đặc trưng trong dữ liệu sinh học.

Trong đồ án này, thuật toán nhóm lựa chọn là FP-Max, một thuật toán khai thác \textit{maximal frequent itemsets} (tập phổ biến tối đại) \cite{Grahne2003FPMax}. FP-Max không đứng riêng lẻ: thuật toán dựa trên FP-tree và các thao tác conditional pattern base của FP-Growth để nén dữ liệu, dựng cây điều kiện và tỉa nhánh. Vì vậy, nhóm phải cài đặt FP-tree/FP-Growth như phần nền tảng bắt buộc trước khi cài FP-Max. Nói cách khác, FP-Growth trong báo cáo này là prerequisite và engine trung gian; FP-Max mới là thuật toán chính của đề tài.

SPMF trong báo cáo không phải là một thuật toán nhóm chọn, mà là một thư viện Java mã nguồn mở về khai thác mẫu (Sequential Pattern Mining Framework) \cite{FournierViger2014SPMF}. Đề bài yêu cầu so sánh với SPMF vì thư viện này có các cài đặt tham chiếu phổ biến cho nhiều thuật toán FIM. Nhóm dùng SPMF như mốc đối chiếu output và thời gian chạy: nếu kết quả của phần FP-tree/FP-Growth nền khớp SPMF, ta có thêm cơ sở để tin rằng các support và cây điều kiện dùng trong FP-Max được xây đúng.

Báo cáo được tổ chức như sau. Mục \ref{sec:theory} trình bày các định nghĩa hình thức của FIM, FP-tree/FP-Growth như nền tảng của FP-Max, giả mã FP-Max và phân tích độ phức tạp. Mục \ref{sec:example} minh hoạ quá trình xây FP-tree rồi rút ra maximal itemsets. Mục \ref{sec:implementation} mô tả cài đặt Julia, trong đó FP-Growth là module nền và FP-Max là module khai thác chính. Mục \ref{sec:evaluation} trình bày kiểm thử, đối chiếu SPMF cho tầng frequent-itemset engine và đánh giá hiệu năng. Mục \ref{sec:application} áp dụng phần frequent-itemset engine vào phân tích giỏ hàng trên tập Groceries để sinh luật kết hợp.
